\documentclass{article}
\usepackage{ACRA/acra}
\usepackage{graphicx}
\usepackage{amsmath}
\usepackage{amssymb}
\usepackage{tabularx}
\usepackage{url}

\title{Context-Conditioned Object Detection for Multi-Robot Perception}
\author{Ysobel Sims \\ Macquarie University, Australia \\ 
ysobel.sims@mq.edu.au}

\begin{document}

\maketitle

\begin{abstract}
Abstract
\end{abstract}

\section{Introduction}

Object detectors typically map an image to a fixed set of semantic categories defined during training. This formulation assumes that the set of objects relevant to detection is determined by the object categories themselves. In robotic environments, however, whether an observed object is relevant can depend on the current state of the environment or task. This state may change during operation without changing the underlying object categories in the scene. A robot may therefore need to apply different detection criteria to the same visual environment at different times.

This distinction can be illustrated by team-based environments. Consider a robot that must detect members of its own team. The robots in the environment remain instances of the same object category, while the visual attribute that determines their relevance depends on the current team configuration. If the team's jersey colour changes, the robot should change which objects it considers relevant without requiring a change to the underlying object detector. More generally, a robot may receive task or environmental information from another system component, a human instruction, a mission planner, or its own perception system. Such information describes the current operating context and can change independently of the objects being detected.

A straightforward solution is to detect all candidate objects and determine their relevance using a separate classification or filtering stage. This can be computationally inexpensive, but requires the contextual attribute to be extracted independently from the detector's learned representation. Attributes affected by object pose, illumination, occlusion, background, or spatial distribution may be difficult to classify reliably using simple post-processing. A separate learned model can provide a more capable contextual classifier, but introduces an additional inference pipeline and associated computational cost. These trade-offs are particularly relevant for small robotic platforms, where perception systems operate under constrained computational resources and compete with other onboard processes.

We therefore investigate incorporating environmental context directly into the detector while retaining a lightweight inference pathway. Rather than treating the contextual state as an additional object category, we provide it as an inference-time conditioning signal that determines which instances of an object category are relevant. This allows the detection criterion to change with the environment while reusing the visual features already computed by the detector. A lightweight conditioning mechanism can consequently provide a middle ground between independent post-processing and an additional learned perception network.

RoboCup Humanoid Soccer provides a controlled setting for investigating this problem. Humanoid robots from opposing teams share the same object category, while team assignment is determined by the current jersey configuration. A robot must therefore distinguish between visually similar instances according to information that can change with the current competition configuration. Jersey colour provides a continuously varying representation of this environmental state: the same detector should identify robots matching the supplied team colour, regardless of which colour is currently assigned to the robot's team. This makes RoboCup a concrete example of a broader robotic setting in which object relevance depends on changing environmental context.

We implement inference-time environmental conditioning using Feature-wise Linear Modulation (FiLM), which modulates the detector's visual features according to a supplied context representation. FiLM provides a compact mechanism for incorporating the conditioning signal into the detector without introducing a separate image-processing network. We evaluate this formulation in a simulated RoboCup Humanoid Soccer environment containing robots with varying jersey colours, comparing direct conditioning with lightweight post-processing approaches, alternative context representations, conditioning mechanisms, and detector architectures.

Our results show that a simple continuous colour representation combined with FiLM provides effective context-conditioned detection while avoiding the additional inference stage required by a separate contextual classifier. In particular, the evaluated post-processing approaches do not reliably recover the contextual attribute, while FiLM provides substantially stronger context-dependent detection with a comparatively lightweight conditioning mechanism.

Our contributions are:

\begin{itemize}
    \item An inference-time conditioning formulation for object detection in which
    externally specified environmental state determines the visual attributes
    relevant to detection.
    
    \item A lightweight FiLM-based implementation that incorporates this
    conditioning into a conventional object detector without introducing
    additional object categories for different environmental configurations.
    
    \item An empirical evaluation in a RoboCup Humanoid Soccer setting, including
    comparisons of context representations, conditioning mechanisms, detector
    architectures, and post-processing approaches.
\end{itemize}

\section{Related Work}


\section{Methods}

\subsection{Dataset}

We generate data using NUpbr \cite{sims2024nupbr}, a physically-based rendering pipeline that produces simulated RoboCup Humanoid Soccer scenes with per-instance ground-truth annotations. For this work, NUpbr's robot-jersey system was extended to support colour-conditioned role assignment: the torso of every robot in a frame is tinted with a jersey colour re-rolled each frame, restricted to at most two distinct colours across all robots in the scene regardless of how many robots are present -- each robot is randomly assigned one of the two colours sampled fresh for that frame. Colours are sampled in HSV space with hue drawn uniformly across the full spectrum and saturation and value each drawn from $[0.6, 1.0]$, producing vividly-coloured, visually distinct jerseys. Each robot's jersey colour is recorded as a \texttt{\#rrggbb} hex string in that robot's entry in the per-frame metadata, alongside its pixel-space bounding box. A scene is retried with a fresh configuration (robot positions, camera target, HDRI environment, and jersey colours) unless between one and four robots are visible in the resulting frame, and unless every visible robot and ball satisfies the annotation gate (in front of the camera, unoccluded, and not implausibly sized due to fisheye edge artifacts); if a robot or ball would otherwise be visible without a valid box, the whole frame is discarded and regenerated rather than emitted with an incomplete annotation.

We additionally post-process the rendered output to remove robots whose jersey is unlikely to be visible: any robot bounding box beginning within a small margin of the top of the frame is dropped before the two-colour and robot-count constraints are (re-)checked, since a box cropped at the top edge typically corresponds to a robot whose torso -- and hence jersey -- lies outside the frame. Images for which more than two distinct jersey colours remain after this filtering are excluded, guaranteeing that every retained image is consistent with a binary teammate/opponent partition of its visible robots.

The dataset used in this work combines two NUpbr render runs, totalling 1,707 images and 2,702 annotated robot instances. Images are partitioned into training, validation, and test splits (70\%/15\%/15\%) at the image level, so that all robots in a given image remain in the same split; since jersey colour is re-sampled independently every frame, there is no instance-level identity to leak across the split boundary. This yields 1,195 training images (1,898 robot boxes), 256 validation images (416 boxes), and 256 test images (388 boxes).

\subsection{Two-Stage Detection and Role Classification}

An early single-stage formulation, in which the jersey-colour embedding directly modulated the detector's backbone features (as in Section~\ref{sec:conditioning}), failed to learn a usable conditioning signal: swapping the supplied teammate and opponent colours at inference time, without retraining, left predictions almost unchanged, indicating that box regression and role classification were competing for the same conditioned representation rather than the model learning to use colour to distinguish roles.

We instead decouple detection from role classification. An unconditioned, single-class ``robot'' detector performs box regression and objectness/classification exactly as it would for any conventional detection task, with no exposure to the jersey-colour embedding. A separate role head then takes an intermediate feature map from the detector, RoIAligns it at each box location -- ground-truth boxes during training, the detector's own predicted boxes at inference -- and predicts teammate or opponent from the resulting box-local feature, conditioned on the supplied jersey-colour embedding. Because the detection pathway never observes the conditioning signal, box quality cannot be corrupted by a difficult-to-learn contextual attribute, and the two components can be evaluated independently.

\subsection{Context Representation}

Each jersey colour is converted from its \texttt{\#rrggbb} hex representation to an RGB triple in $[0,1]^3$. The conditioning signal supplied to the model is the concatenation of the teammate and opponent colours, a six-dimensional embedding $(r_t, g_t, b_t, r_o, g_o, b_o)$. During training, the assignment of which of an image's (up to two) jersey colours is ``teammate'' and which is ``opponent'' is re-randomised on every access to the training set, so that the same image is seen under both role assignments across epochs; this prevents the model from learning shortcuts tied to a fixed colour-to-role mapping rather than to the supplied conditioning signal. Validation and test sets instead use one fixed, randomly drawn role assignment per image, so that evaluation metrics remain comparable across epochs and runs.

\subsection{Conditioning Mechanism}
\label{sec:conditioning}

The primary conditioning mechanism is Feature-wise Linear Modulation (FiLM): a small generator network maps the six-dimensional colour embedding to a per-channel scale and shift, which are applied to the RoIAligned feature before pooling and classification. We additionally evaluate a cross-attention mechanism, in which spatial positions of the pooled feature attend to the embedding (treated as a single token), as an alternative to FiLM's affine modulation.

\subsection{Detector Backbones}

We evaluate three detector backbones: Faster R-CNN with a ResNet-50-FPN backbone, YOLO26n, and RT-DETR. Each is pretrained on a general-purpose object detection dataset before fine-tuning. Since the three architectures expose intermediate feature maps at different points in their computational graphs, the role head's RoIAlign tap is chosen separately for each: an FPN level for Faster R-CNN, an early backbone-level feature for YOLO26n, and the analogous pre-encoder-fusion backbone feature for RT-DETR, matched by resolution and position in the network to keep the comparison as controlled as possible across backbones.

\subsection{Non-Learned and Decoupled Baselines}

To establish how much of the role head's accuracy is attributable to learned visual representations rather than colour alone, we evaluate a set of non-learned post-processing rules that classify a box's role by nearest Euclidean RGB distance to the two candidate jersey colours, using either the mean colour of the box (or a central crop of it) or the minimum distance over a grid of cells within the box.

We also evaluate a standalone classifier: a dedicated, lightweight convolutional encoder with its own FiLM-conditioned classification head, trained independently of any detector directly on ground-truth boxes. Unlike the integrated role head, this classifier does not share features with a detection backbone and instead operates on a RoIAligned crop of the raw image, so it can in principle be deployed behind any box-producing detector without requiring access to that detector's internal representations.

\subsection{Training Details}

Images are letterboxed to $640\times640$ and augmented with random horizontal flipping (probability 0.5) during training. Models are optimised with AdamW, with learning rate and batch size tuned per architecture. Training uses early stopping with patience evaluated jointly across detection loss, role loss, and role accuracy (or, for the detector-only and standalone-classifier ablations, the subset of these that apply), so that training does not halt while any one of the jointly-tracked metrics is still improving.

\section{Experiments}

We report role classification accuracy, computed by applying the trained role head to ground-truth boxes and comparing against the known jersey colour of each box, and detection accuracy (map$_{50}$), computed from the detector's own predicted boxes, for all detector-based configurations.

\textbf{Backbone comparison.} We compare Faster R-CNN, RT-DETR, and YOLO26n, each pretrained and paired with a FiLM-only role head (no auxiliary colour-distance feature), across three random seeds (Table~\ref{tab:backbone-comparison}).

\textbf{Role-head architecture comparison.} Holding the Faster R-CNN backbone fixed, we compare three role-head variants: FiLM conditioning combined with an auxiliary hand-crafted colour-distance feature, FiLM conditioning alone, and cross-attention conditioning, again across three seeds (Table~\ref{tab:role-head-comparison}).

\textbf{Role classification approaches.} We compare the non-learned colour-distance rules, the standalone classifier, and the integrated FiLM role head, all evaluated on the same ground-truth test boxes, to isolate the contribution of learned representations and of sharing features with a detection backbone from the contribution of colour information alone (Table~\ref{tab:role-classification}).

\textbf{Swap consistency.} To test whether the role head uses the supplied colour embedding causally rather than exploiting some other correlate of the image, we evaluate each trained model twice on the same ground-truth boxes: once with the true teammate/opponent embedding, and once with the two colours swapped. A causally-conditioned model should flip its prediction under a swap whenever it was correct under the true embedding; we report this as the conditional probability that a swapped prediction is correct (relative to the flipped label) given that the corresponding unswapped prediction was correct (Table~\ref{tab:swap-consistency}).

\section{Results}

% Table 1: Backbone comparison
\begin{table}[h]
\centering
\begin{tabular}{l|cc}
\textbf{Backbone} & \textbf{Role Acc (\%)} & \textbf{map$_{50}$ (\%)} \\
\hline
Faster R-CNN & $92.1 \pm 1.5$ & $84.7 \pm 2.8$ \\
RT-DETR & $82.7 \pm 5.3$ & $69.0 \pm 7.6$ \\
YOLO26n & $79.2 \pm 2.8$ & $61.5 \pm 4.6$ \\
\end{tabular}
\caption{Backbone comparison, all with a pretrained backbone and a FiLM-only (RGB) role head, mean and standard deviation over 3 seeds. Note: FiLM-only is the stronger role head on Faster R-CNN and RT-DETR, but the \emph{weaker} one on YOLO26n, where FiLM+distance instead scores $82.6 \pm 0.6$ / $64.6 \pm 1.2$ -- held fixed here for a controlled comparison across backbones.}
\label{tab:backbone-comparison}
\end{table}

% Table 2: Role-head architecture comparison
\begin{table}[h]
\centering
\begin{tabular}{l|cc}
\textbf{Role Head} & \textbf{Role Acc (\%)} & \textbf{map$_{50}$ (\%)} \\
\hline
FiLM (RGB) + dist & $90.1 \pm 1.8$ & $81.0 \pm 1.0$ \\
\hline
FiLM (RGB) & $92.1 \pm 1.5$ & $84.7 \pm 2.8$ \\
\hline
Cross-attention & $83.5 \pm 1.8$ & $70.2 \pm 3.7$ \\
\end{tabular}
\caption{Role-head architecture comparison, all on a pretrained Faster R-CNN (ResNet-50) backbone, mean $\pm$ std over 3 seeds.}
\label{tab:role-head-comparison}
\end{table}


% Table 3: Comparison of role classification approaches
\begin{table}[t]
\centering
\small
\begin{tabular}{l|c|c}
\textbf{Method} & \textbf{Role Acc. (\%)} & \textbf{Parameters} \\
\hline
Mean RGB, whole box & 48.7 & -- \\
Min RGB distance, whole box & 50.3 & -- \\
Mean RGB, 15\% centre crop & 50.5 & -- \\
Min RGB distance, 15\% centre crop & 54.4 & -- \\
\hline
Standalone classifier (RGB, learned) & $87.1 \pm 2.3$ & 161,538 \\
FiLM (RGB, learned, integrated) & $\mathbf{92.1 \pm 1.5}$ & \mathbf{133,890} \\
\end{tabular}
\caption{Comparison of role classification approaches on ground-truth boxes ($n=388$). Non-learned methods assign each box to the nearest of the two injected reference colours using Euclidean RGB distance; ``mean RGB'' uses the mean box/crop colour, while ``min RGB distance'' assigns the box using the minimum distance over cells. The standalone classifier is a dedicated lightweight CNN with a FiLM-conditioned classification head, trained independently of the detector. The integrated FiLM head instead operates on features already computed by the detection backbone. Learned results are mean $\pm$ std over 3 seeds using a pretrained Faster R-CNN. Parameter counts exclude the shared detection backbone. The standalone classifier comprises a 93,696-parameter dedicated CNN encoder and a 67,842-parameter FiLM conditioner and classifier, whereas the integrated FiLM head adds 133,890 parameters to the detector.}
\label{tab:role-classification}
\end{table}

The standalone classifier achieves substantially higher role accuracy than the non-learned colour rules, confirming that role assignment benefits from learned visual representations. However, its dedicated encoder adds 161,538 parameters to the system. The integrated FiLM head achieves higher accuracy while adding only 133,890 parameters, since it operates directly on features already computed for detection.

\begin{table}[h]
\centering
\begin{tabular}{l|c|c}
Metric & YOLO26n & Faster R-CNN \\
\hline
Normal accuracy & 74.7 & 89.7 \\
\hline
Swap accuracy (vs.\ flipped label) & 72.2 & 91.2 \\
\hline
Both correct / dataset & 65.2 & 86.6 \\
\hline
P(swap correct $\mid$ normal correct) & 87.2 & 96.6 \\
\end{tabular}
\caption{Swap consistency on ground-truth boxes (test split, $n=388$): the role head is run once with the true teammate/opponent RGB embedding and once with it swapped, for both FiLM-only backbones. All values are percentages.}
\label{tab:swap-consistency}
\end{table}



\section{Conclusion} 

% wildlife: the same species might have different types of calls - aggressive, mating, etc

\section*{Acknowledgments}


\bibliographystyle{ACRA/named}
\bibliography{X1-citations}

\end{document}

